\chapter{Analyse de l'existant}

\section{Présentation de la solution actuelle}

Actuellement, l'ANTIC ne dispose pas d'un outil centralisé dédié à la conversion automatisée de fichiers Excel en documents Word. Les agents utilisent les outils bureautiques standards installés sur leurs postes de travail individuels (Microsoft Office ou LibreOffice), combinés le cas échéant à des scripts VBA locaux ou à des convertisseurs en ligne non agréés.

\section{Fonctionnalités et fonctionnement général actuel}

Le processus actuel d'élaboration manuelle des rapports Word à partir de tableaux de bord tableurs est illustré par le diagramme d'activités ci-dessous (Figure~\ref{fig:diagramme_activite_existant}) :

\begin{figure}[H]
\centering
\begin{tikzpicture}[
    node distance=0.9cm,
    startend/.style={draw=primarycolor, fill=primarycolor!15, rectangle, rounded corners=12pt, minimum width=3.8cm, minimum height=0.8cm, align=center, font=\sffamily\bfseries\small},
    activity/.style={draw=secondarycolor, fill=secondarycolor!10, rectangle, rounded corners=5pt, minimum width=6.2cm, minimum height=0.85cm, align=center, font=\sffamily\small, inner sep=5pt},
    decision/.style={draw=orange!80!black, fill=orange!10, diamond, aspect=2, minimum width=3.6cm, minimum height=1.0cm, align=center, font=\sffamily\small, inner sep=2pt},
    arrow/.style={draw=primarycolor, ->, >=stealth, thick}
]

    % Nodes
    \node (start) [startend] {Début du processus};
    \node (step1) [activity, below=of start] {1. Ouverture du classeur tableur (Excel / Calc)};
    \node (step2) [activity, below=of step1] {2. Sélection \& filtrage manuel des données};
    \node (step3) [activity, below=of step2] {3. Transfert manuel par Copier-Coller vers Word};
    \node (step4) [activity, below=of step3] {4. Réajustement manuel des styles et colonnes};
    \node (dec) [decision, below=of step4] {Tous les onglets\\traités ?};
    \node (end) [startend, below=of dec] {Rapport Word finalisé};

    % Main flow arrows
    \draw [arrow] (start) -- (step1);
    \draw [arrow] (step1) -- (step2);
    \draw [arrow] (step2) -- (step3);
    \draw [arrow] (step3) -- (step4);
    \draw [arrow] (step4) -- (dec);
    \draw [arrow] (dec) -- node[right=4pt, font=\sffamily\small\bfseries\color{primarycolor}] {Oui} (end);
    
    % Loop back arrow for 'Non' (wide clear margin on left)
    \draw [arrow] (dec.west) -- node[above=2pt, font=\sffamily\small\bfseries\color{orange!80!black}] {Non} ++(-3.2,0) |- (step2.west);

\end{tikzpicture}
\caption{Diagramme d'activités du processus manuel actuel}
\label{fig:diagramme_activite_existant}
\end{figure}

\section{Acteurs du système actuel}

Les acteurs intervenant dans ce processus manuel sont l'\textbf{ensemble des employés de l'ANTIC} (toutes directions et tous services confondus) impliqués dans la manipulation des données tableurs et la rédaction des rapports officiels.

\section{Analyse de la première approche (\textit{SheetTools})}

Une première expérimentation interne (Proof of Concept - POC) via l'Add-In Excel \textbf{SheetTools.xlam} visait à automatiser le filtrage dynamique et l'exportation vers Word au moyen de macros VBA locales.

Bien qu'elle ait permis de valider les principes métier (création de sous-feuilles et mise en page Paysage / AutoFit), cette approche locale n'a \textbf{pas été retenue pour un déploiement officiel}. Son abandon s'explique par sa dépendance aux postes clients Windows, l'instabilité du presse-papier local, les risques sécuritaires liés aux macros et l'absence d'audit centralisé.
